\documentclass[../main.tex]{subfiles}

\begin{document}

\ifSubfilesClassLoaded{

    %\tableofcontents
    
    %\setcounter{chapter}{}
}{}

\chapter{ Important Distribution and Their Calculations }
% 代靖涵 25120222201319
\section{Discrete Distributions}
\begin{proposition}{Bernoulli Distribution}{}
Let $p \in [0,1]$. If a random variable $X$ takes values in $\{0,1\}$ and satisfies $P(X=1)=p, P(X=0)=1-p$, then $X$ is said to follow a \dfntxt{Bernoulli distribution} with parameter $p$, denoted as $X \sim \text{Bernoulli}(p)$. In this case:

\begin{enumerate}
\item The probability mass function is 
\[P(X=k) = p^k(1-p)^{1-k}, \quad k \in \{0,1\};\]

\item $E[X] = p$;

\item $\text{Var}(X) = p(1-p)$;

\item The moment generating function is 
\[M_X(t) = 1-p+pe^t.\]
\end{enumerate}
\end{proposition}

\begin{proposition}{Binomial Distribution}{}
Let $n \in \mathbb{N}^+, p \in [0,1]$. If a random variable $X$ represents the number of successes in $n$ independent repeated Bernoulli trials, then $X$ is said to follow a \dfntxt{binomial distribution} with parameters $(n,p)$, denoted as $X \sim B(n,p)$. In this case:

\begin{enumerate}
\item The probability mass function is 
\[P(X=k) = \binom{n}{k}p^k(1-p)^{n-k}, \quad k=0,1,\ldots,n;\]

\item $E[X] = np$;

\item $\text{Var}(X) = np(1-p)$;

\item The moment generating function is 
\[M_X(t) = (1-p+pe^t)^n;\]

\item If $X \sim B(n_1,p), Y \sim B(n_2,p)$ and $X,Y$ are independent, then 
\[X+Y \sim B(n_1+n_2,p).\]
\end{enumerate}
\end{proposition}

\begin{proposition}{Geometric Distribution}{}
Let $p \in (0,1]$. If a random variable $X$ represents the number of trials needed to obtain the first success in independent repeated Bernoulli trials, then $X$ is said to follow a \dfntxt{geometric distribution} with parameter $p$. In this case:

\begin{enumerate}
\item The probability mass function is 
\[P(X=k) = (1-p)^{k-1}p, \quad k=1,2,3,\ldots;\]

\item \[E[X] = \frac{1}{p};\]

\item \[\text{Var}(X) = \frac{1-p}{p^2};\]

\item Memoryless property: for any $m,n \in \mathbb{N}^+$, we have 
\[P(X>m+n|X>m) = P(X>n);\]

\item The moment generating function is 
\[M_X(t) = \frac{pe^t}{1-(1-p)e^t}, \quad t < -\ln(1-p).\]
\end{enumerate}
\end{proposition}
\begin{proof}
     \[
     \begin{aligned}
     P\left( X> m+ n|X> m \right) = \frac{P\left( X> m+ n \right)  }{P\left( X> m \right)  }= \frac{\sum _{i= m+ n+ 1}^{\infty}\left( 1-p \right)^{i-1}p  }{ \sum _{i= m+ 1}^{\infty}\left( 1-p \right)^{i-1}p }  &= \frac{\frac{\left( 1-p \right)^{m+ n}  }{p }  }{\frac{\left( 1-p \right)^{m}  }{ p}  }  \\ 
      &= \frac{\left( 1-p \right)^{n}  }{p }p\\ 
       &=  \sum _{i= n+ 1}^{\infty}\left( 1-p \right)^{i-1}p\\ 
        &= P\left( X> n \right)  
     \end{aligned}
     \]
\end{proof}
\begin{proposition}{Negative Binomial Distribution}{}
Let $r \in \mathbb{N}^+, p \in (0,1]$. If a random variable $X$ represents the number of trials needed to obtain the $r$-th success in independent repeated Bernoulli trials, then $X$ is said to follow a \dfntxt{negative binomial distribution} with parameters $(r,p)$, denoted as $X \sim NB(r,p)$. In this case:

\begin{enumerate}
\item The probability mass function is 
\[P(X=k) = \binom{k-1}{r-1}p^r(1-p)^{k-r}, \quad k=r,r+1,\ldots;\]

\item \[E[X] = \frac{r}{p};\]

\item \[\text{Var}(X) = \frac{r(1-p)}{p^2};\]

\item When $r=1$, the negative binomial distribution reduces to the geometric distribution;

\item If $X \sim NB(r_1,p), Y \sim NB(r_2,p)$ and $X,Y$ are independent, then 
\[X+Y \sim NB(r_1+r_2,p).\]
\end{enumerate}
\end{proposition}

\begin{proposition}{Poisson Distribution}{}
Let $\lambda > 0$. If a random variable $X$ takes values in $\mathbb{N} \cup \{0\}$ with probability mass function 
\[P(X=k) = \frac{\lambda^k e^{-\lambda}}{k!},\]
then $X$ is said to follow a \dfntxt{Poisson distribution} with parameter $\lambda$, denoted as $X \sim \text{Poisson}(\lambda)$. In this case:

\begin{enumerate}
\item $E[X] = \lambda$;

\item $\text{Var}(X) = \lambda$;

\item The moment generating function is 
\[M_X(t) = e^{\lambda(e^t-1)};\]

\item If $X \sim \text{Poisson}(\lambda_1), Y \sim \text{Poisson}(\lambda_2)$ and $X,Y$ are independent, then 
\[X+Y \sim \text{Poisson}(\lambda_1+\lambda_2).\]
\end{enumerate}
\end{proposition}

\begin{proposition}{Poisson Limit Theorem}{}
Let $X_n \sim B(n,p_n)$, where $np_n \to \lambda > 0$ as $n \to \infty$. Then for any $k \in \mathbb{N} \cup \{0\}$, we have
\[\lim_{n \to \infty} P(X_n = k) = \frac{\lambda^k e^{-\lambda}}{k!}.\]
\end{proposition}

\begin{proposition}{Hypergeometric Distribution}{}
Let $N,M,n \in \mathbb{N}^+$ with $M \leq N, n \leq N$. If a random variable $X$ represents the number of success elements when drawing $n$ elements without replacement from a population containing $M$ success elements and $N-M$ failure elements, then $X$ is said to follow a \dfntxt{hypergeometric distribution} with parameters $(N,M,n)$, denoted as $X \sim H(N,M,n)$. In this case:

\begin{enumerate}
\item The probability mass function is 
\[P(X=k) = \frac{\binom{M}{k}\binom{N-M}{n-k}}{\binom{N}{n}}, \quad k=\max(0,n-N+M),\ldots,\min(n,M);\]

\item \[E[X] = n\frac{M}{N};\]

\item \[\text{Var}(X) = n\frac{M}{N}\cdot\frac{N-M}{N}\cdot\frac{N-n}{N-1};\]

\item As $N \to \infty$ and $\frac{M}{N} \to p$, the hypergeometric distribution converges in distribution to the binomial distribution $B(n,p)$.
\end{enumerate}
\end{proposition}
\section{Continuous Distribution}
\begin{proposition}{Uniform Distribution}{均匀分布}
Let $a, b \in \mathbb{R}$ with $a < b$. If a random variable $X$ has a probability density function (PDF) that is constant over the interval $[a,b]$, then $X$ is said to follow a \dfntxt{uniform distribution} on $[a,b]$, denoted as $X \sim U(a,b)$. In this case:
\begin{enumerate}
\item The probability density function is
$$\begin{cases} 
\frac{1}{b-a} & \text{if } a \le x \le b, \\
0 & \text{otherwise}.
\end{cases}$$
\item $E[X] = \frac{a+b}{2}$;
\item $\text{Var}(X) = \frac{(b-a)^2}{12}$;
\item The moment generating function is
$M_X(t) = \frac{e^{tb}-e^{ta}}{t(b-a)}, \quad \text{for } t \neq 0.$
\end{enumerate}
\end{proposition}
\begin{proposition}{Exponential Distribution}{指数分布}
Let $\lambda > 0$. If a random variable $X$ represents the waiting time for an event to occur in a Poisson process, it is said to follow an \dfntxt{exponential distribution} with rate parameter $\lambda$, denoted as $X \sim \text{Exp}(\lambda)$. In this case:
\begin{enumerate}
\item The probability density function is
$$\begin{cases} 
\lambda e^{-\lambda x} & \text{if } x \ge 0, \\
0 & \text{otherwise}.
\end{cases}$$
\item $E[X] = \frac{1}{\lambda}$;
\item $\text{Var}(X) = \frac{1}{\lambda^2}$;
\item Memoryless property: for any $s, t \ge 0$, we have
$P(X > s+t | X > s) = P(X > t);$
\item The moment generating function is
$M_X(t) = \frac{\lambda}{\lambda-t}, \quad t < \lambda.$
\end{enumerate}
\end{proposition}
\begin{proposition}{Normal Distribution}{正态分布}
Let $\mu \in \mathbb{R}$ and $\sigma^2 > 0$. A random variable $X$ is said to follow a \dfntxt{normal distribution} (or \dfntxt{Gaussian distribution}) with parameters $(\mu, \sigma^2)$, denoted as $X \sim N(\mu, \sigma^2)$. In this case:
\begin{enumerate}
\item The probability density function is
$f_X(x) = \frac{1}{\sigma\sqrt{2\pi}} e^{-\frac{(x-\mu)^2}{2\sigma^2}}, \quad x \in \mathbb{R};$
\item $E[X] = \mu$;
\item $\text{Var}(X) = \sigma^2$;
\item The moment generating function is
$M_X(t) = e^{\mu t + \frac{1}{2}\sigma^2 t^2};$
\item If $X \sim N(\mu, \sigma^2)$, then the standardized variable $Z = \frac{X-\mu}{\sigma}$ follows the \dfntxt{standard normal distribution} $N(0,1)$;
\item If $X \sim N(\mu_1, \sigma_1^2)$, $Y \sim N(\mu_2, \sigma_2^2)$ and $X,Y$ are independent, then for any constants $a,b \in \mathbb{R}$,
$aX+bY \sim N(a\mu_1+b\mu_2, a^2\sigma_1^2+b^2\sigma_2^2).$
\end{enumerate}
\end{proposition}
\begin{proposition}{Lognormal Distribution}{对数正态分布}
Let $Y$ be a positive random variable. If the new random variable $X = \ln(Y)$ follows a normal distribution with parameters $(\mu, \sigma^2)$, then $Y$ is said to follow a \dfntxt{lognormal distribution}, denoted as $Y \sim LN(\mu, \sigma^2)$. Note that $\mu$ and $\sigma^2$ are the mean and variance of the logarithm of the variable, not of the variable itself.
\begin{enumerate}
\item The probability density function is
$$f_Y(y) = 
\begin{cases} 
\frac{1}{y\sigma\sqrt{2\pi}} \exp\left(-\frac{(\ln y - \mu)^2}{2\sigma^2}\right) & \text{if } y > 0, \\
0 & \text{otherwise}.
\end{cases}$$
\item $E[Y] = e^{\mu + \frac{1}{2}\sigma^2}$;
\item $\text{Var}(Y) = (e^{\sigma^2}-1)e^{2\mu + \sigma^2}$;
\item The key relationship for problem-solving: If $Y \sim LN(\mu, \sigma^2)$, then $\ln(Y) \sim N(\mu, \sigma^2)$. This allows us to transform problems about lognormal variables into problems about normal variables.
\item The lognormal distribution arises from the multiplicative version of the Central Limit Theorem. The product of many independent, identically distributed positive random variables will be approximately log-normally distributed.
\end{enumerate}
\end{proposition}
\begin{proposition}{Gamma Distribution}{伽玛分布}
Let $\alpha > 0$ (shape parameter) and $\lambda > 0$ (rate parameter). A random variable $X$ is said to follow a \dfntxt{Gamma distribution} with parameters $(\alpha, \lambda)$, denoted as $X \sim \Gamma(\alpha, \lambda)$. In this case:
\begin{enumerate}
\item The probability density function is
$$\begin{cases} 
\frac{\lambda^\alpha}{\Gamma(\alpha)} x^{\alpha-1} e^{-\lambda x} & \text{if } x > 0, \\
0 & \text{otherwise},
\end{cases}$$
where $\Gamma(\alpha) = \int_0^\infty z^{\alpha-1}e^{-z}dz$ is the Gamma function.
\item $E[X] = \frac{\alpha}{\lambda}$;
\item $\text{Var}(X) = \frac{\alpha}{\lambda^2}$;
\item The moment generating function is
$M_X(t) = \left(\frac{\lambda}{\lambda-t}\right)^\alpha, \quad t < \lambda;$
\item If $X \sim \Gamma(\alpha_1, \lambda)$, $Y \sim \Gamma(\alpha_2, \lambda)$ and $X,Y$ are independent, then
$X+Y \sim \Gamma(\alpha_1+\alpha_2, \lambda);$
\item The exponential distribution is a special case: $\text{Exp}(\lambda) \equiv \Gamma(1, \lambda)$;
\item The \dfntxt{Chi-squared distribution} (卡方分布) with $k$ degrees of freedom, denoted $\chi^2_k$, is a special case: $\chi^2_k \equiv \Gamma(k/2, 1/2)$.
\end{enumerate}
\end{proposition}
\begin{proposition}{Central Limit Theorem}{中心极限定理}
Let $X_1, X_2, \ldots$ be a sequence of independent and identically distributed (i.i.d.) random variables with finite mean $E[X_i] = \mu$ and finite variance $\text{Var}(X_i) = \sigma^2 > 0$. Let $S_n = \sum_{i=1}^n X_i$. Then the standardized sample mean converges in distribution to a standard normal distribution:
$\frac{\bar{X}_n - \mu}{\sigma/\sqrt{n}} = \frac{S_n - n\mu}{\sigma\sqrt{n}} \xrightarrow{d} N(0,1) \quad \text{as } n \to \infty.$
\end{proposition}

\begin{proposition}{Beta Distribution}{贝塔分布}
Let $\alpha > 0$ and $\beta > 0$ be shape parameters. A random variable $X$ is said to follow a \dfntxt{Beta distribution} with parameters $(\alpha, \beta)$, denoted as $X \sim \text{Beta}(\alpha, \beta)$, if its support is the interval $(0,1)$. This distribution is commonly used to model probabilities or proportions.
\begin{enumerate}
\item The probability density function is
$$f_X(x) = 
\begin{cases} 
\frac{x^{\alpha-1}(1-x)^{\beta-1}}{B(\alpha, \beta)} & \text{if } 0 < x < 1, \\
0 & \text{otherwise},
\end{cases}$$
where $B(\alpha, \beta) = \frac{\Gamma(\alpha)\Gamma(\beta)}{\Gamma(\alpha+\beta)}$ is the Beta function.
\item $E[X] = \frac{\alpha}{\alpha+\beta}$;
\item $\text{Var}(X) = \frac{\alpha\beta}{(\alpha+\beta)^2(\alpha+\beta+1)}$;
\item The $k$-th moment is $E[X^k] = \frac{\Gamma(\alpha+k)\Gamma(\alpha+\beta)}{\Gamma(\alpha)\Gamma(\alpha+\beta+k)}$ for any $k > -\alpha$;
\item The moment generating function does not have a simple closed form.
\item If $U \sim \Gamma(\alpha, \lambda)$ and $V \sim \Gamma(\beta, \lambda)$ are independent, then the random variable $X = \frac{U}{U+V}$ follows a Beta distribution, $X \sim \text{Beta}(\alpha, \beta)$.
\end{enumerate}
\end{proposition}

\begin{note}
    Beta分布是在我们相信一个观测结果是由二项分布产生的情况下, 导致这个观测结果的 \(  p  \)的分布, \(  p  \)是我们认为的二项分布的参数.  
\end{note}

\begin{proposition}{Cauchy Distribution}{柯西分布}
Let $x_0 \in \mathbb{R}$ be the location parameter (median) and $\gamma > 0$ be the scale parameter. A random variable $X$ is said to follow a \dfntxt{Cauchy distribution}, denoted as $X \sim \text{Cauchy}(x_0, \gamma)$. It is a classic example of a distribution for which the moments (like mean and variance) are undefined. The case with $x_0=0, \gamma=1$ is called the \dfntxt{standard Cauchy distribution}.
\begin{enumerate}
\item The probability density function is
$$f_X(x) = \frac{1}{\pi\gamma \left[1 + \left(\frac{x-x_0}{\gamma}\right)^2\right]}, \quad x \in \mathbb{R};$$
\item $E[X]$ is undefined;
\item $\text{Var}(X)$ is undefined;
\item The moment generating function is undefined. However, its characteristic function is well-defined: $\phi_X(t) = E[e^{itX}] = e^{ix_0 t - \gamma|t|}$;
\item If $U, V \sim N(0,1)$ are independent standard normal variables, then their ratio $U/V$ follows the standard Cauchy distribution, $\text{Cauchy}(0,1)$;
\item The Cauchy distribution is a special case of the Student's t-distribution with one degree of freedom.
\item If $X_1, \ldots, X_n$ are i.i.d. from $\text{Cauchy}(x_0, \gamma)$, their sample mean $\bar{X}_n$ has the same distribution, $\bar{X}_n \sim \text{Cauchy}(x_0, \gamma)$. This shows it does not obey the Law of Large Numbers.
\end{enumerate}
\end{proposition}
\begin{proposition}{Student's t-Distribution}{学生t-分布}
Let $\nu > 0$ be the degrees of freedom. A random variable $T$ is said to follow a \dfntxt{Student's t-distribution} with $\nu$ degrees of freedom, denoted as $T \sim t_\nu$. This distribution arises when estimating the mean of a normally distributed population in situations where the sample size is small and the population standard deviation is unknown.
\begin{enumerate}
\item The probability density function is
$$f_T(t) = \frac{\Gamma((\nu+1)/2)}{\sqrt{\nu\pi}\Gamma(\nu/2)} \left(1 + \frac{t^2}{\nu}\right)^{-\frac{\nu+1}{2}}, \quad t \in \mathbb{R};$$
\item $E[T] = 0$ for $\nu > 1$. It is undefined for $\nu \le 1$;
\item $\text{Var}(T) = \frac{\nu}{\nu-2}$ for $\nu > 2$. It is infinite for $1 < \nu \le 2$ and undefined for $\nu \le 1$;
\item If $Z \sim N(0,1)$ and $V \sim \chi^2_\nu$ are independent, then $T = \frac{Z}{\sqrt{V/\nu}} \sim t_\nu$;
\item As $\nu \to \infty$, the t-distribution converges in distribution to the standard normal distribution $N(0,1)$.
\end{enumerate}
\end{proposition}
\begin{proposition}{F-Distribution}{F-分布}
Let $d_1 > 0$ and $d_2 > 0$ be the numerator and denominator degrees of freedom, respectively. A random variable $X$ is said to follow an \dfntxt{F-distribution}, denoted as $X \sim F(d_1, d_2)$. It is most commonly used in the analysis of variance (ANOVA).
\begin{enumerate}
\item The probability density function is
$$f_X(x) = 
\begin{cases} 
\frac{1}{B(d_1/2, d_2/2)} \left(\frac{d_1}{d_2}\right)^{d_1/2} x^{d_1/2 - 1} \left(1 + \frac{d_1}{d_2}x\right)^{-(d_1+d_2)/2} & \text{if } x > 0, \\
0 & \text{otherwise};
\end{cases}$$
\item $E[X] = \frac{d_2}{d_2-2}$ for $d_2 > 2$;
\item $\text{Var}(X) = \frac{2d_2^2(d_1+d_2-2)}{d_1(d_2-2)^2(d_2-4)}$ for $d_2 > 4$;
\item If $U_1 \sim \chi^2_{d_1}$ and $U_2 \sim \chi^2_{d_2}$ are independent, then the ratio of the scaled variables $X = \frac{U_1/d_1}{U_2/d_2}$ follows an F-distribution, $X \sim F(d_1, d_2)$;
\item If $T \sim t_{d_2}$, then $T^2 \sim F(1, d_2)$.
\end{enumerate}
\end{proposition}

\section{Mutivariable Distributions}

\begin{proposition}{Multivariate Normal Distribution}{多元正态分布}
Let $\boldsymbol{\mu} \in \mathbb{R}^k$ be a mean vector and $\boldsymbol{\Sigma}$ be a $k \times k$ symmetric, positive-definite covariance matrix. A random vector $\mathbf{X} = (X_1, \ldots, X_k)^T$ is said to follow a \dfntxt{multivariate normal distribution}, denoted as $\mathbf{X} \sim N_k(\boldsymbol{\mu}, \boldsymbol{\Sigma})$.
\begin{enumerate}
\item The probability density function is
$f_{\mathbf{X}}(\mathbf{x}) = \frac{1}{(2\pi)^{k/2} |\boldsymbol{\Sigma}|^{1/2}} \exp\left(-\frac{1}{2}(\mathbf{x}-\boldsymbol{\mu})^T \boldsymbol{\Sigma}^{-1}(\mathbf{x}-\boldsymbol{\mu})\right), \quad \mathbf{x} \in \mathbb{R}^k;$
where $|\boldsymbol{\Sigma}|$ is the determinant of $\boldsymbol{\Sigma}$. The term $(\mathbf{x}-\boldsymbol{\mu})^T \boldsymbol{\Sigma}^{-1}(\mathbf{x}-\boldsymbol{\mu})$ is the Mahalanobis distance squared.
\item $E[\mathbf{X}] = \boldsymbol{\mu}$;
\item $\text{Cov}(\mathbf{X}) = E[(\mathbf{X}-\boldsymbol{\mu})(\mathbf{X}-\boldsymbol{\mu})^T] = \boldsymbol{\Sigma}$;
\item The moment generating function is
$M_{\mathbf{X}}(\mathbf{t}) = E[e^{\mathbf{t}^T\mathbf{X}}] = e^{\mathbf{t}^T\boldsymbol{\mu} + \frac{1}{2}\mathbf{t}^T\boldsymbol{\Sigma}\mathbf{t}};$
\item \textbf{Marginal Distributions}: Any subset of the components of $\mathbf{X}$ also follows a multivariate normal distribution.
\item \textbf{Linear Transformations}: If $\mathbf{Y} = \mathbf{A}\mathbf{X} + \mathbf{b}$ for a matrix $\mathbf{A}$ of size $m \times k$ and vector $\mathbf{b} \in \mathbb{R}^m$, then $\mathbf{Y} \sim N_m(\mathbf{A}\boldsymbol{\mu}+\mathbf{b}, \mathbf{A}\boldsymbol{\Sigma}\mathbf{A}^T)$.
\item Two components $X_i$ and $X_j$ are independent if and only if their covariance $\boldsymbol{\Sigma}_{ij} = 0$. This is a special and powerful property of the normal distribution.
\end{enumerate}
\end{proposition}
\begin{proposition}{Multinomial Distribution}{多项分布}
Let $n \in \mathbb{N}^+$ be the number of trials, and $\mathbf{p} = (p_1, \ldots, p_k)$ be a vector of probabilities such that $p_i \ge 0$ for all $i$ and $\sum_{i=1}^k p_i = 1$. If a random vector $\mathbf{X} = (X_1, \ldots, X_k)$ represents the number of times each of the $k$ possible outcomes occurs in $n$ independent trials, then $\mathbf{X}$ follows a \dfntxt{multinomial distribution}, denoted as $\mathbf{X} \sim \text{Mult}(n, \mathbf{p})$.
\begin{enumerate}
\item The probability mass function is
$P(X_1=n_1, \ldots, X_k=n_k) = \frac{n!}{n_1! n_2! \cdots n_k!} p_1^{n_1} p_2^{n_2} \cdots p_k^{n_k},$
where $n_i \in \{0, 1, \ldots, n\}$ and $\sum_{i=1}^k n_i = n$.
\item $E[X_i] = np_i$;
\item $\text{Var}(X_i) = np_i(1-p_i)$;
\item $\text{Cov}(X_i, X_j) = -np_i p_j$ for $i \neq j$. The negative covariance reflects the constraint that $\sum X_i = n$.
\item The binomial distribution is a special case where $k=2$.
\end{enumerate}
\end{proposition}
\begin{proposition}{Dirichlet Distribution}{狄利克雷分布}
Let $\boldsymbol{\alpha} = (\alpha_1, \ldots, \alpha_k)$ be a vector of concentration parameters with $\alpha_i > 0$. A random vector $\mathbf{X} = (X_1, \ldots, X_k)$ is said to follow a \dfntxt{Dirichlet distribution}, denoted as $\mathbf{X} \sim \text{Dir}(\boldsymbol{\alpha})$. Its support is the standard $(k-1)$-simplex, i.e., $S_{k-1} = \{\mathbf{x} \in \mathbb{R}^k : x_i > 0, \sum_{i=1}^k x_i = 1\}$. It is the multivariate generalization of the Beta distribution.
\begin{enumerate}
\item The probability density function is
$f(\mathbf{x}; \boldsymbol{\alpha}) = \frac{1}{B(\boldsymbol{\alpha})} \prod_{i=1}^k x_i^{\alpha_i - 1},$
where $B(\boldsymbol{\alpha}) = \frac{\prod_{i=1}^k \Gamma(\alpha_i)}{\Gamma(\sum_{i=1}^k \alpha_i)}$ is the multivariate Beta function.
\item Let $\alpha_0 = \sum_{i=1}^k \alpha_i$.
\item $E[X_i] = \frac{\alpha_i}{\alpha_0}$;
\item $\text{Var}(X_i) = \frac{\alpha_i(\alpha_0 - \alpha_i)}{\alpha_0^2(\alpha_0+1)}$;
\item $\text{Cov}(X_i, X_j) = \frac{-\alpha_i \alpha_j}{\alpha_0^2(\alpha_0+1)}$ for $i \neq j$;
\item The Dirichlet distribution is the conjugate prior for the parameter vector $\mathbf{p}$ of the multinomial distribution.
\end{enumerate}
\end{proposition}
\begin{proposition}{Wishart Distribution}{威沙特分布}
Let $\mathbf{V}$ be a $p \times p$ symmetric positive-definite scale matrix and let $n \ge p$ be the degrees of freedom. A random matrix $\mathbf{W}$ is said to follow a \dfntxt{Wishart distribution}, denoted $\mathbf{W} \sim W_p(n, \mathbf{V})$, if it is the distribution of the sample covariance matrix. Specifically, if $\mathbf{X}_1, \ldots, \mathbf{X}_n$ are i.i.d. samples from $N_p(\mathbf{0}, \mathbf{V})$, then $\mathbf{W} = \sum_{i=1}^n \mathbf{X}_i \mathbf{X}_i^T \sim W_p(n, \mathbf{V})$.
\begin{enumerate}
\item The probability density function for a symmetric positive-definite matrix $\mathbf{X}$ is
$f(\mathbf{X}) = \frac{|\mathbf{X}|^{(n-p-1)/2} \exp\left(-\frac{1}{2}\text{tr}(\mathbf{V}^{-1}\mathbf{X})\right)}{2^{np/2} |\mathbf{V}|^{n/2} \Gamma_p(n/2)},$
where $\Gamma_p(\cdot)$ is the multivariate Gamma function.
\item $E[\mathbf{W}] = n\mathbf{V}$;
\item The Wishart distribution is the multivariate generalization of the Chi-squared distribution. If $p=1$ and $\mathbf{V}=1$, then $W_1(n,1)$ is the $\chi^2_n$ distribution.
\item \textbf{Additive Property}: If $\mathbf{W}_1 \sim W_p(n_1, \mathbf{V})$ and $\mathbf{W}_2 \sim W_p(n_2, \mathbf{V})$ are independent, then $\mathbf{W}_1 + \mathbf{W}_2 \sim W_p(n_1+n_2, \mathbf{V})$.
\end{enumerate}
\end{proposition}
\begin{proposition}{Uniform Distribution on the Sphere}{球面上均匀分布}
This distribution describes a point chosen at random from the surface of a unit sphere in $\mathbb{R}^n$. Let $S^{n-1} = \{\mathbf{x} \in \mathbb{R}^n : ||\mathbf{x}||_2 = 1\}$ be the unit sphere. A random vector $\mathbf{X}$ is uniformly distributed on $S^{n-1}$. This is a foundational example of a distribution on a Riemannian manifold.
\begin{enumerate}
\item The probability density function is constant with respect to the surface area measure (the canonical volume form on the sphere). If $A_{n-1}$ is the surface area of $S^{n-1}$, the density is $f(\mathbf{x}) = 1/A_{n-1}$.
\item The surface area of $S^{n-1}$ is given by $A_{n-1} = \frac{2\pi^{n/2}}{\Gamma(n/2)}$.
\item \textbf{Generation}: If $\mathbf{Z} = (Z_1, \ldots, Z_n)^T$ where $Z_i \sim N(0,1)$ are i.i.d., then the vector $\mathbf{X} = \frac{\mathbf{Z}}{||\mathbf{Z}||_2}$ is uniformly distributed on $S^{n-1}$. This is a consequence of the rotational symmetry of the multivariate normal distribution.
\item $E[\mathbf{X}] = \mathbf{0}$ due to symmetry.
\item $\text{Cov}(\mathbf{X}) = E[\mathbf{X}\mathbf{X}^T] = \frac{1}{n}\mathbf{I}_n$, where $\mathbf{I}_n$ is the $n \times n$ identity matrix.
\end{enumerate}
\end{proposition}

\begin{theorem}{}{}
    使得
\end{theorem}

\end{document}